\scanpage{101}
\mspage{96}
\begin{center}
Chapter III

Some Applications of the Criterion of Projectivity
of Abstract Complete Varieties.
\end{center}

In this chapter we will examine some of the applications
of the following assertion which follows from theorem 2:

\uline{Let} $\mathrm{V}$ \uline{be a nonsingular complete algebraic variety}
$\mathrm{\dim V=s}$ \uline{and} $\mathrm{D}$ \uline{be a divisor on} $\mathrm{V}$ \uline{such that} $\mathrm{D^k W_k>0}$
\uline{for any irreducible subvariety} $\mathrm{W_k}$ \uline{of dimension}
$\mathrm{k}$, $\mathrm{k=\overline{1,s}}$ $\mathrm{(D^sV=D^s)}$. \uline{Then, at sufficiently large} $\mathrm{n}$,
\uline{the mapping of the variety} $\mathrm{V}$ \uline{into a projective space},
\uline{corresponding to the fibre} $\mathrm{[nD]}$, \uline{becomes a biregular
imbedding}.

Let $\mathrm{V}$ be a complete normal algebraic variety,
$\mathrm{\dim V=s}$ and $\mathrm{U}$ be an open set on $\mathrm{V}$, being a quasi-projective
variety (that is, allowing realization in the
\margin{p. 218}form of an open subset of a projective algebraic variety).
Let $\mathrm{x_o,x_1,....}$ be homogeneous coordinates, corresponding
to $\mathrm{U}$, $\mathrm{x_o\ne 0}$, and $\mathrm{F_j=\frac{x_j}{x_o}}$. Let $\mathrm{\{F\}}$ be \deleted{a}\inserted{\handwritten{\unsure{the}}} linear space
over $\mathrm{k}$, generated by the functions $\mathrm{F_j}$, and $\mathrm{\{T_{\deleted{\unsure{1|l}}\handwritten{\ell}}\}}$ be a
\margin{\handwritten{$\mathrm{\ell}$}}
totality of irreducible $\mathrm{s-1}$ dimensional components of
the set $\mathrm{V-U}$,
\[
\handwritten{\tau_\ell} = \handwritten{\min_j n_{T_\ell}(F_j)}
\]
where $\mathrm{n_{T_{\unsure{l|1}}}(F_j)}$ denotes the order of the function $\mathrm{F_j}$ on
the divisor $\mathrm{T_l}$. Let a function $\mathrm{\unsure{\varphi|\phi}\in\{F\}}$ be such that,
% archive sheet number: 100

